\documentclass[10pt]{exam} % Usa la clase 'exam'

% Configuración de la cabecera
\usepackage[utf8]{inputenc}
\usepackage[T1]{fontenc}
\usepackage{amsmath} % Por si necesitas matemáticas

\title{Decentralized Certification on Bitcoin Blockchain Using Taproot}
\author{Student: Juan~Crespo~Vargas\\Advisor: Dr.~Sultan~Ahmad~Al-Muhammadi}
\date{}
\begin{document}

\maketitle


\begin{abstract}
    
Decentralized digital certification offers an alternative to traditional centralized trust models by reducing dependency on intermediary servers and improving institutional autonomy and system availability.
However, current blockchain-based solutions such as Blockcerts or smart contract platforms on Ethereum, face limitations in consensus security, issuer privacy, and operational cost sustainability.
The absence of a proof-of-work mechanism in these networks increases the risk of chain reorganizations and denial-of-service attacks, undermining the immutability required for long-term certificate verification.

This research proposes a digital certification system built on Bitcoin’s Layer 1, leveraging the capabilities introduced by Taproot (BIP-341/342) and the Schnorr digital signature scheme to ensure authenticity, proof of existence, and verifiable key rotation for certificate issuers.
By combining Merkle trees and off-chain verification proofs, the system aims to achieve scalability and privacy without overloading the blockchain with public data.
The proposal also includes a mitigation scheme against key loss or compromise through cryptographic policies embedded within Taproot scripts.
Finally, the study evaluates implementation costs, computational complexity, and privacy guarantees in comparison with existing standards such as OpenTimestamps and Blockcerts, to assess the technical feasibility and contribution of the proposed model to decentralized trust infrastructure.

\end{abstract}
\end{document}
